% !TEX root = ../main.tex
\chapter{Discussion}
\label{ch:discussion}

Chapter~\ref{ch:results} reported the measurements taken on Arbitrum One. Here each research question and each objective of Chapter~\ref{ch:introduction} is held to the criterion stated there, and every interpretation points to the table of Chapter~\ref{ch:results} it rests on. The measurements are those of the revised analysis, which values every amount in US dollars and keeps five tokens. The amounts were revised after the registered results were known (Section~\ref{sec:amount-revision}), so none of the revised results discussed here is a registered result, and a registered result, where one is mentioned, is labelled as such; the registered analysis, in raw base units over every token, reached the same verdicts (Appendix~\ref{sec:map-registered}, Table~\ref{tab:registered-vs-revised}). Unless noted otherwise, same-window statements concern the matched cohort of contracts and holdout statements the spender cohort.

\dissertationSubheading[sec:findings-rq]{Findings by research question}

\dissertationSubsubheading[sub:finding-rq1]{RQ1: Reputation structure and rank divergence}

RQ1 asks whether EndorseRank and AWP order the same contracts in different ways. They do: the interval of their agreement excludes both zero and one (Table~\ref{tab:tie-sensitivity}), and of the contract pairs that both scores order, 41.2\% are ordered in opposite directions (Section~\ref{sec:rank-divergence}). Ties play almost no part in this result. Every cohort contract received approvals of one of the five tokens from at least three owners, so EndorseRank leaves few pairs of contracts tied, and keeping the contracts that are missing from a graph as isolated nodes, instead of scoring them zero, leaves the agreement unchanged.

The divergence follows from how the two graphs are built. The allowance graph records the latest approval of each owner, spender and token, whereas the transfer graph records every transfer \parencite{do2023awp}, and the allowance graph has fewer than half of the transfer graph's edges (Section~\ref{sec:benchmark-results}); the two methods also restart differently. PageRank places stationary mass according to the edges it is given \parencite{page1999pagerank}, so a contract that receives many or large transfers from active senders but few approvals ranks high under AWP and modestly under EndorseRank, and the reverse holds for a contract that many owners approve without being a transfer hub. H1 is supported, largely by construction, since two walks over different graphs with different restarts were not expected to agree. In practice one score cannot stand in for the other: an operator who adopts EndorseRank because its graph is smaller gets a different ranking and not a cheaper approximation of AWP.

\dissertationSubsubheading[sub:finding-rq2]{RQ2: Same-window alignment}

RQ2 asks how closely each score agrees with contemporaneous allowance and transfer checks. Each score follows the edges it walks on (Table~\ref{tab:alignment-hybrid}). EndorseRank and the allowance layer walked alone agree most with the allowance family, AWP and the transfer layer walked alone most with the transfer family, and the paired contrasts between the two single-layer scores exclude zero: EndorseRank agrees more with the allowance family than with the transfer family, and more with the allowance family than AWP does, while AWP agrees more with the transfer family than EndorseRank does (Table~\ref{tab:tau-diff}). The C-PR mixtures lie between their two layers on every family and move far less between $\lambda=0.25$ and $\lambda=0.75$ than between the two layers. Every interval in Table~\ref{tab:alignment-hybrid} excludes zero except EndorseRank's on the Sybil-stability family, so by the criterion of O2 EndorseRank does not agree with that construct, and every other score agrees with each construct to some degree; what separates the scores is the family they agree with most.

Two caveats bound H2. The first is circularity. Every proxy in a family is computed from the same events as the matching graph, so the agreement is an intended-construct check in the sense of \textcite{cronbach1955construct}, and part of it is built in: the allowance proxy with which EndorseRank agrees most is the in-approve degree, the count it smooths (Table~\ref{tab:alignment-allowance}), and AWP, whose edges weigh each transfer by its dollar value, follows the summed value of inbound transfers almost as closely as the number of its senders (Table~\ref{tab:alignment-transfer}). Ties do not distort the comparison, since EndorseRank leaves few pairs tied and the isolated-node rule changes no coefficient (Table~\ref{tab:tie-sensitivity}). The second caveat is that a higher coefficient in the same window does not mark a better score out of window. EndorseRank with AWP's activity restarts agrees less with the allowance family than EndorseRank does ($\tau=0.426$, $[0.418,0.435]$, against $0.585$, $[0.578,0.592]$; no paired contrast was computed, and the two intervals do not overlap), yet it ranks both spender labels better in both holdout windows, with paired intervals above zero, in comparisons made after the labels were known (Table~\ref{tab:endorserank-restarts} and Section~\ref{sub:endorserank-restarts}). In the narrow sense of agreement with the checks of a score's own edges, H2 is supported, largely by construction. The paired differences between EndorseRank and AWP, which carry no claim, are in Appendix~\ref{ch:appendix-er-awp}.

\dissertationSubsubheading[sub:finding-rq3]{RQ3: Computational cost}

RQ3 asks what each score costs to compute on the same sample. The cost follows the edge count. On the same contracts AWP's graph has $2.23$ times the edges of EndorseRank's, and its solver call takes $2.17$ times as long (Section~\ref{sec:benchmark-results}; Table~\ref{tab:benchmark-runtime} gives $|V|$, $|E|$, peak memory and iterations of each graph). C-PR and S-PR also walk the transfer layer and take longer than AWP. Matrix assembly and iteration are both linear in $|E|$ \parencite{langville2006pagerank}, so the time ratio is expected to track the edge ratio. Both scores run the same solver, so the lower cost is a property of the sparser graph and not of the algorithm. Each graph and setting was timed once under the protocol of five timed runs after one warm-up, and that measurement is reused wherever the graph appears.

The scaling stages repeat the comparison on smaller graphs. Each stage keeps the edges incident to a deterministic subsample of the cohort, so both edge counts grow from stage to stage (Table~\ref{tab:benchmark-scaling}). The time ratio of the two solves follows the edge ratio from the stage of $10{,}000$ contracts on, but not at the three smaller stages, where the time ratio departs from the edge ratio in either direction (Section~\ref{sub:runtime-scaling}). The stages are subsets of one ego network, so they say nothing about either solver on a chain-wide graph. The damping and value-slope checks leave the order of the families as it was and move the family means little (Tables~\ref{tab:robustness-damping} and~\ref{tab:robustness-slope}); with the slopes of the value weight divided or multiplied by ten, each score still agrees most with its own family and ranks the contracts almost as the main score does. H3 is supported on the matched cohort and its subsamples, again largely by construction, since both scores run the same solver. H1--H3 are thus descriptive checks of the instruments; the tests whose outcome was open are RQ4, RQ5 and Section~\ref{sub:finding-neutral}.

\dissertationSubsubheading[sub:finding-rq4]{RQ4: Temporal holdout of future new approvals}

RQ4 has two clauses. The first asks whether scores frozen at the end of March 2026 predict new owner--spender approval pairs and new transfer senders on the spender cohort over the following three months. Every score predicts both labels, and the edge type decides how well (Table~\ref{tab:holdout-spenders}). Among the scores of that table, the interior mixtures and the in-approve degree lead on new approval pairs, followed by EndorseRank and the allowance layer, while the scores built on transfers alone stay lower. On new senders AWP leads, ahead of the transfer in-degree and the transfer layer, and EndorseRank, the allowance layer and the in-approve degree rank lowest. The window W1, with scores frozen at the end of June and labels counted from July to September, repeats this pattern among the same scores, except that on new approval pairs the in-approve degree has a higher point estimate than the interior mixtures, a difference for which no paired contrast was computed (Tables~\ref{tab:fresh-holdout} and~\ref{tab:fresh-posthoc}). The evidence concerns spender contracts; Section~\ref{sub:finding-neutral} turns to the traders.

The second clause asks whether any PageRank beats the raw degree it smooths, and the two walks answer it differently (Tables~\ref{tab:holdout-diff} and~\ref{tab:fresh-posthoc}). AWP ranks later new senders better than the transfer in-degree in both windows, with intervals above zero, and the lead holds on the spenders that had received a transfer before the freeze, so the spenders that AWP scores zero do not create it (Table~\ref{tab:holdout-receiving}). On the transfer side, then, the walk with its time and value weights and activity restarts added information that the count lacks. EndorseRank, with the uniform restarts fixed in the analysis plan, ranks new approval pairs below the in-approve degree in both windows, and so does the allowance layer walked alone in W0. On the allowance side the count was the better predictor. Known results lead one to expect the count to be a hard baseline to beat. In growing networks new links attach preferentially to nodes that already have many \parencite{barabasi1999scaling}, which makes degree a standard link-prediction baseline \parencite{libennowell2007link}, also on Ethereum transaction networks \parencite{lin2022tracking}. Where degree correlations are weak, as on the Web, PageRank is on average proportional to in-degree \parencite{fortunato2008indegree}. For most spenders EndorseRank is an approval count in which each approval is weighted by its age and dollar amount and divided by the summed weight of all the approvals its owner grants (Section~\ref{sub:theory-closed-form}), which leaves it little room to differ from the count.

The variant of EndorseRank with AWP's activity restarts ranks new approval pairs above the in-approve degree in both windows (Section~\ref{sub:endorserank-restarts} and Table~\ref{tab:fresh-posthoc}). Only the restart rule separates the two variants, so on the allowance side the restart rule appears to decide whether the PageRank beats its count; whether the same holds for AWP was not tested, since AWP was run only with its published restarts. The variant's contrasts were computed after the labels were known in both windows, so they state a hypothesis for a future registered test and not a finding. None of the degree contrasts was fixed in advance: on the revised scores every one of them was computed after the labels of its window were known. In the registered analysis the four degree contrasts of EndorseRank and AWP had the same signs, with intervals on the same side of zero (Table~\ref{tab:registered-vs-revised}). They do not undo the first clause, since the counts are statistics of the same edges.

The labels lie outside the score window but belong to the same constructs as the scores, and neither contains any element of profit, loss or repayment. RQ4 therefore shows how far each construct carries forward in time, with the raw degree as the baseline that a PageRank must beat, and leaves the safety and the credit risk of a contract aside; the liquidation label of W1 (Section~\ref{sub:fresh-holdout}), on which Section~\ref{sub:finding-neutral} compares the edge types, comes closest to a risk outcome. The revocation counts and the drain heuristic computed on the same spenders rise with every score, and the revocations most with the in-approve degree (Section~\ref{sec:holdout-results}). They track the volume of approvals a contract holds, and they are not offered as evidence about its safety.

\dissertationSubsubheading[sub:finding-rq5]{RQ5: Coupling the two edge types}

Both halves of RQ5 get a negative answer, each under its own rule. The rules were fixed for the registered analysis and are applied here to the revised scores after the registered results were known, so the verdicts below are a re-analysis and not registered tests; the registered analysis reached the same verdicts (Table~\ref{tab:registered-vs-revised}). The second half asks whether one walk over both layers ranks contracts better than either layer walked alone. Under rule A, fixed in the analysis plan before any score existed, it does not. On W0, C-PR is no worse than the allowance layer on new approval pairs but falls behind the transfer layer on new senders (Table~\ref{tab:holdout-diff}), and in the same window its point estimates lie outside the interval of the better layer on both the transfer and the allowance family and between the two layers on Sybil stability (Section~\ref{sec:coupled-results}), so the secondary part fails on all three families as well. The first half asks whether the walk improves on the transfer layer walked alone. An exploratory comparison on W0 said that it did on new approval pairs, but it was made after the labels were known. Rule B, registered before the July--September logs were extracted, was written for this question, and on the revised scores it is not met. C-PR again beat the transfer layer on new approval pairs, by a little more than in W0, but it fell behind on new senders by more than the margin of $0.01$, and on the traders' liquidation-free close rate it could not show that it was no worse: the two point estimates are equal, but the interval of their difference reaches below $-0.01$ (Section~\ref{sub:fresh-results}). On that trader label every score has an interval above zero (Table~\ref{tab:fresh-holdout}). Against AWP, the comparator of the first sensitivity analysis, C-PR falls behind on new senders by a wider margin. Against EndorseRank, in comparisons not fixed in advance, it is slightly above on new approval pairs in both windows (Tables~\ref{tab:holdout-diff} and~\ref{tab:fresh-posthoc}).

The ablation and the layer weights explain both halves. S-PR keeps the transfer walk but restarts it from the allowance walk. In contrasts computed after the labels were known, it stays close to the transfer layer on new senders and falls far below the allowance layer on new approval pairs (Table~\ref{tab:holdout-diff}), since a teleportation vector derived from endorsements changes little when transfer edges decide where the walk goes. In C-PR the allowance edges carry part of the walk, so the allowance signal survives on the spender holdout: every interior mixture has a slightly higher coefficient on new approval pairs than the allowance layer alone, and at $\lambda=0.5$ the paired interval lies above zero (Tables~\ref{tab:holdout-spenders} and~\ref{tab:holdout-diff}). The price is paid on the transfer side, where every interior mixture has a lower coefficient on new senders than the transfer layer. C-PR's layers carry dollar amounts without a bound, apart from the cap on allowances, and restart uniformly, whereas AWP, which bounds the weight of each transfer through $V$ and restarts at active senders, ranks new senders better than C-PR in both windows (after the labels in W0, under the first sensitivity analysis applied to the revised scores in W1). On these data allowance edges help a walk only when they carry it. Other couplings were not tested (Section~\ref{sec:limitations}).

Both halves count as results. Rule A was fixed before any coupled score existed, so $\lambda$ could not have been tuned to it; the revision changed the tokens, the cohort and the amounts in its layers, adding the cap $U$ to the allowance layer, after its registered verdict was known, but neither $\lambda$ nor its thresholds. It limits an obvious next step: a convex mixture of the two edge types at the transition level gained less over the allowance layer on new approval pairs than it lost against the transfer layer on new senders, and in the same window it stayed outside the better layer's interval on both families (Tables~\ref{tab:holdout-diff} and~\ref{tab:alignment-hybrid}). Rule B was registered before its labels were extracted, and on the revised scores, as in the registered analysis, it fails on two of its three conditions. W1 shows something narrower: adding allowance edges to a transfer walk gains information about authorization that holds on new data, at a loss on inflow larger than the margin fixed in the registration.

\dissertationSubsubheading[sub:finding-neutral]{Allowance against transfer on outcomes built from neither edge type}

Unlike the spender labels of RQ4 and RQ5, which favor the scores built on their own edge type, the GMX trader outcomes are built from neither edge type and were measured after the scores were frozen. The comparison of the two edge types on them (Section~\ref{sub:neutral-label-protocol}) adds no research question and is reported as Claim~C6 (Section~\ref{sub:neutral-label-results}). It was registered together with rule B, before the July--September logs were extracted and before the W0 liquidation labels had been computed for any score, and the registration fixed the wording of its conclusions. Here its rules R1--R3 and their wording are applied to the revised scores, after the registered results were known, so the verdicts are a re-analysis and not registered tests; the registered analysis reached the same three verdicts (Table~\ref{tab:registered-vs-revised}).

This comparison bears on counterparty risk, the use that motivated the thesis (Chapter~\ref{ch:introduction}). For a protocol that extends leverage or credit to a contract account, a counterparty that is often force-liquidated is a riskier one, and forced liquidation is the nearest outcome to default in these data. In the primary window W1 the paired difference between the approval count at the freeze and the transfer count on the liquidation-free close rate is $\Delta\tau=+0.084$ (95\% interval $[-0.013,0.176]$, Bonferroni 97.5\% interval $[-0.026,0.188]$), and stratified by the number of closes it is $+0.101$ ($[-0.024,0.231]$); both intervals include zero. In W0 the same contrast was $+0.182$ ($[0.052,0.314]$), but W0 is the secondary window. The like-for-like PageRank pair, EndorseRank with AWP's restarts against AWP, is negative in both windows with intervals that include zero (Table~\ref{tab:neutral-label-contrasts}). Under the wording fixed in the registration, the result therefore reads: no count-level lead of allowance over transfer on later liquidation avoidance; the PageRank versions are not resolved; and, since the approval count ranks traders lower than the transfer count on the profitable share and on the non-loss share, the signal concerns forced-liquidation avoidance, not trading success.

What remains is descriptive. Every trader that had received an approval by the freeze avoided liquidation in both windows, 37 of 37 in W1 and 40 of 40 in W0, against about three in four of the other traders (Table~\ref{tab:neutral-label-recipients}). Among the traders the in-approve degree is at most one in W1 and three in W0, so the approval count works here as a flag for having been approved at all, and the transfer count also ranks the label in W1 (Table~\ref{tab:fresh-holdout}).

Three limits bound this reading. First, the samples are small. Few contract accounts close positions on GMX~V2 in these months, so the trader cohorts hold 70 and 88 contracts and every interval is wide. Second, account type cannot explain the pattern, because every trader in both cohorts is a contract, but the role of a contract may: the recipients close more positions than the other traders and were authorized by other addresses to spend their tokens, a pattern that fits managed or pooled trading contracts, which the data do not identify. Third, the recipients have lower profitable and non-loss shares, so what they avoid is forced liquidation, not loss; the lower non-loss share is partly built in, because the label counts a liquidated close as not a loss. Within these limits, the comparison did not separate allowance-edge from transfer-edge information on the outcome that a margin or lending protocol screens for.

\dissertationSubheading[sec:objectives-achievement]{Achievement of the research objectives}

Chapter~\ref{ch:introduction} attached a measurement and a criterion to each of its five objectives. Table~\ref{tab:objectives-achievement} checks them against the evidence of the revised analysis. Where a criterion is a rule fixed in the plan or the registration, the outcome is that of the rule applied to the revised scores after the registered results were known; the registered outcomes are in Table~\ref{tab:registered-vs-revised}.

\begingroup
\setstretch{1}\small\setlength{\tabcolsep}{3.5pt}\captionsetup{font=normalsize}
\begin{longtable}{p{0.05\linewidth}p{0.33\linewidth}p{0.44\linewidth}p{0.12\linewidth}}
\caption{Research objectives, the criterion stated in Chapter~\ref{ch:introduction}, and the outcome.}
\label{tab:objectives-achievement}\\
\toprule
Obj. & Criterion & Evidence & Outcome \\
\midrule
\endfirsthead
\caption[]{(continued)}\\
\toprule
Obj. & Criterion & Evidence & Outcome \\
\midrule
\endhead
\bottomrule
\endlastfoot
O1 & 95\% interval of inter-method $\tau$ contains neither $0$ nor $1$ & Table~\ref{tab:tie-sensitivity}, last row: $\tau=0.174$, interval $[0.163,0.185]$, the same under the isolated-node rule & Met (largely by construction) \\
O2 & A score agrees with a construct when its interval on that family excludes $0$; agreement with the construct of a score's own edges is an intended-construct check & Table~\ref{tab:alignment-hybrid}: every interval excludes $0$ except EndorseRank's on Sybil stability ($0.007$, $[-0.001,0.015]$); each single-layer score agrees most with its own family (EndorseRank allowance $0.585$, $[0.578,0.592]$; AWP transfer $0.572$, $[0.564,0.580]$), and the paired contrasts of Table~\ref{tab:tau-diff} exclude $0$; the proxies come from the same events as the graphs & Met (largely by construction) \\
O3 & Time ratio reported only together with the ratio of $|E|$, under one fixed protocol & Section~\ref{sec:benchmark-results} and Table~\ref{tab:benchmark-runtime}: five timed runs after one warm-up, with $|V|$ and $|E|$ of each solver graph; AWP's time ratio to EndorseRank $2.17$ beside an edge ratio of $2.23$; one solver, so time follows $|E|$ & Met (largely by construction) \\
O4 & A score predicts a label when its interval excludes $0$; a PageRank improves on its degree only when the interval of the paired difference lies above $0$ & Tables~\ref{tab:holdout-spenders}, \ref{tab:fresh-holdout} and~\ref{tab:fresh-posthoc}: every score predicts both labels in W0 and W1. Tables~\ref{tab:holdout-diff} and~\ref{tab:fresh-posthoc}: AWP minus transfer in-degree $+0.043$ $[0.038,0.048]$ (W0) and $+0.022$ $[0.017,0.027]$ (W1); EndorseRank minus in-approve degree $-0.019$ $[-0.026,-0.011]$ (W0) and $-0.034$ $[-0.038,-0.029]$ (W1); C-PR ($\lambda=1$) minus in-approve degree $-0.027$ $[-0.034,-0.019]$ (W0); all computed after the labels were known & Partly met (first clause met; second clause met for AWP, not for EndorseRank) \\
O5 & (a) Rule A (analysis plan, 9~October~2026): on W0, C-PR does no worse than C-PR ($\lambda=1$) on new approval pairs and better than C-PR ($\lambda=0$) on new senders; in the same window, its point estimate lies within the better layer's interval on the transfer and allowance families, and it exceeds both layers on Sybil stability. (b) Rule B (registered at 16:49~UTC on 9~October~2026, before the July--September logs were extracted): on W1, C-PR beats C-PR ($\lambda=0$) on new approval pairs (F1) and is not worse by more than $\delta=0.01$ on new senders (F2) and on the traders' liquidation-free close rate (F3) & (a) Table~\ref{tab:holdout-diff}: no worse holds ($+0.029$, $[0.022,0.036]$), better fails ($-0.034$, $[-0.040,-0.029]$). Same window (Table~\ref{tab:alignment-hybrid}): transfer $0.365$ $[0.355,0.375]$ outside $[0.653,0.666]$; allowance $0.477$ $[0.469,0.485]$ outside $[0.590,0.604]$; Sybil stability $0.142$ $[0.134,0.151]$, between the layers. (b) Table~\ref{tab:fresh-contrasts}: F1 $+0.110$ $[0.098,0.122]$ holds; F2 $-0.028$ $[-0.033,-0.023]$ and F3 $+0.000$ $[-0.014,0.010]$ fail & (a) Not met; (b) not met; both on the revised scores, as in the registered analysis \\
\end{longtable}
\endgroup

Six findings outside the objectives also shape the conclusions. Three limit the claims. The sample-definition stages move the same-window coefficients enough that they cannot serve as robustness evidence (Table~\ref{tab:robustness-sample-size}); the value-slope check bounds the effect of the slopes chosen in the revision only in the same window, where dividing or multiplying them by ten keeps each score closest to its own family and its ranking close to the main one (Table~\ref{tab:robustness-slope}), and no other slope was tried on the holdout labels; and on the trader cohorts no score ranks the share of profitable closes with an interval above zero in either window (Tables~\ref{tab:holdout-traders} and~\ref{tab:fresh-traders}). Three add to the claims. EndorseRank with AWP's activity restarts ranked new approval pairs above the in-approve degree in both windows, in comparisons made after the labels and offered as a hypothesis (Section~\ref{sub:finding-rq4}). Every trader contract that had received an approval avoided liquidation in both windows, in small samples and without a lead over the transfer count under rule R1 (Section~\ref{sub:finding-neutral}). And EndorseRank matches a closed form to within 1\% for every depth-one spender, while the graph also has deeper paths along which propagation changes the score (Section~\ref{sub:theory-closed-form}). That AWP beats its count and EndorseRank does not is not among them, because it is the second clause of O4.

\dissertationSubheading[sec:theoretical-implications]{Theoretical implications}

The results bear on how propagation works on financial graphs, on what EndorseRank computes for spenders whose owners are not themselves approved, on the construct alignment of on-chain scores and on trust under cheap pseudonyms.

\dissertationSubsubheading[sub:theory-propagation]{Propagation semantics beyond transfers}

PageRank was first formulated for hyperlinks, read as a citation network \parencite{page1999pagerank}. \textcite{do2023hodgerank} carried this reading over to Ethereum transactions, and AWP \parencite{do2023awp} refined it with value and age weights and restarts at active senders. EndorseRank applies it to ERC-20 allowances, where an owner who lets a spender contract act for it issues a directed endorsement of spending authority \parencite{vogelsteller2015erc20}. The RQ1 divergence shows that changing what the edges mean, together with the restart rule that goes with each graph, changes the ranking even on the same nodes, so graph-based reputation is a family of methods, each tied to the relationships its edges encode. Peer-to-peer and web trust systems reached the same conclusion from the other side: in EigenTrust \parencite{kamvar2003eigentrust} and TrustRank \parencite{gyongyi2004trustrank} the seed set and the edge definitions largely decide the outcome. EndorseRank's edge is native to a token standard, where theirs belonged to a message-passing overlay or a hyperlink graph.

The holdouts add a qualification that the propagation metaphor hides. Changing the edge type changed what the score could anticipate: on new approval pairs the scores built on allowance edges led and those built on transfers alone stayed lower, and on new senders the reverse held, in both windows. Whether propagating along the edges added to the count depended on the side. On transfers AWP beat its count; on allowances EndorseRank did not, while its variant with activity restarts did. None of these comparisons was fixed in advance, and on the revised scores all of them were made after the labels were known (Section~\ref{sub:finding-rq4}). On these data the edge type sets what a score anticipates, and the restart rule, more than the walk itself, appears to decide whether the score improves on the count. When the two layers share one walk, the allowance signal survives where allowance edges carry the walk and is lost where they only seed its restarts (Section~\ref{sub:finding-rq5}).

\dissertationSubsubheading[sub:theory-closed-form]{EndorseRank on depth-one spenders}

The allowance graph keeps the approvals of the five tokens that a cohort contract grants or receives (Section~\ref{sec:graph-construction}). Most of its nodes are pure owners, with out-edges only, or pure spenders, with in-edges only. A cohort contract can be both, since it can receive approvals from owners and approve other contracts in turn, so the graph has paths of more than one step. For a spender whose owners receive no approvals, EndorseRank has a closed form.

\emph{Proposition.} Let PageRank be computed by Equation~\eqref{eq:pagerank} with damping $d$ and uniform restarts, so that dangling mass is also spread uniformly. Then every node without in-edges has the same score $c$, and every node $v$ whose in-neighbors all lack in-edges has
\begin{equation}
\label{eq:closed-form}
r_v = c\,\bigl(1 + d\,\delta(v)\bigr), \qquad \delta(v)=\sum_{u} \frac{w(u,v)}{o(u)},
\end{equation}
where $\delta(v)$ is the owner-normalized in-strength of $v$: each owner's endorsement is divided over all the spenders it approves.

\emph{Derivation.} At the fixed point, $r_v = (1-d+d\,r_D)/|V| + d\sum_u r_u\,w(u,v)/o(u)$, where $r_D$ is the mass held by the dangling nodes. A node without in-edges receives only the first term, so every such node has the score $c=(1-d+d\,r_D)/|V|$, the constant that Section~\ref{sub:sybil-bound} also writes $c$. If every in-neighbor $u$ of $v$ is such a node, substituting $r_u=c$ gives Equation~\eqref{eq:closed-form}.

On the graph at $T_{\mathrm{obs}}$ the nodes without in-edges share one score, $c=7.64\times10^{-8}$, to numerical precision. Of the 14{,}406 cohort spenders in that graph, 12{,}039 have depth one, and EndorseRank as computed in this thesis lies within 1\% of Equation~\eqref{eq:closed-form} for every one of them. Of the other 2{,}367, which have an owner that itself receives approvals, 26.4\% lie within 1\%. Over all spenders the share is 87.9\%. Of the 17{,}032 W0 spenders in the graph at $t_1$, 11{,}326 have depth one, every one of them again within 1\%, and the share over all W0 spenders is 73.7\%. The share of deeper spenders is larger there partly because the W0 cohort also contains contracts outside the matched cohort, which enter the extracted data only through approvals from cohort contracts, and every cohort contract receives approvals of its own in the observation window (Section~\ref{sub:cohorts}). The check is computed by \texttt{scripts/run\_same\_window.py} and, for W0, \texttt{scripts/run\_holdout.py} of the experiment code, run with the switch that selects the revised analysis (Section~\ref{sec:reproducibility}).

Two consequences follow. First, for a depth-one spender EndorseRank only reweights the approval count: each owner's latest approval of a token counts by its age and its dollar amount, never for more than one, and an owner that approves many spenders counts for less at each. These edge weights and that normalization are all that separate it from the count, and for these spenders there is no propagated information for the walk to add. Second, the graph does have deeper paths: about a sixth of the cohort spenders at $T_{\mathrm{obs}}$ and a third of the W0 spenders have an owner that is itself endorsed, and for most of them the score departs from the closed form. Multi-hop propagation is therefore present in this graph, but it did not make EndorseRank with uniform restarts beat the raw count out of window (Section~\ref{sub:finding-rq4}). Whether it helps on a deeper sample, such as the two-hop neighborhood of the cohort or a chain-wide allowance graph, or with restarts seeded at vetted contracts, is not tested here.

\dissertationSubsubheading[sub:theory-construct]{Multi-construct domain alignment}

The checks were designed to separate transfer centrality from allowance delegation, and each score aligns with the construct of the edges it walks on, the mixtures falling between their layers. Classical construct-validity theory would predict this pattern for instruments that measure different constructs \parencite{cronbach1955construct}, and \textcite{messick1995validity} goes further, counting the consequences of using a score as part of its validity. This thesis applies that logic to ordinal reputation scores with Spearman's $\rho$ \parencite{spearman1904proof} and Kendall's $\tau$ \parencite{kendall1938rank}, which need no classification threshold, and the paired bootstrap \parencite{efron1979bootstrap} makes the pattern testable. The same theory says what such a pattern cannot settle. When two instruments measure different constructs, the one that comes out ahead on a check is the one built for that check, so ordering the instruments by their coefficients does not show which is the better measure of reputation. The two EndorseRank variants add a sharper point: the variant that agrees less with its own construct in the same window anticipates that construct's later events better, in comparisons made after the labels were known (Section~\ref{sub:finding-rq2}), so same-window agreement does not rank instruments even within one construct. Outcomes built from neither edge type are the exception, since neither instrument was built for them, which is why the thesis compares the edge types there (Section~\ref{sub:finding-neutral}); in the primary window that comparison, its registered rules applied to the revised scores, did not separate them. Such comparisons are possible here because no score leaves many pairs tied; with a heavily tied score, two rank correlations could be compared only when they share the same attainable maximum.

\dissertationSubsubheading[sub:theory-pseudonymity]{Reputation under cheap pseudonyms}

Where identities cost little to create, one entity can present many \parencite{douceur2002sybil}, and under Ethereum's account model a new address needs no registration step at all \parencite{wood2025yellowpaper}. Any reputation system that rewards inbound transfers or approvals is therefore exposed to Sybil strategies, in which an attacker creates edges among addresses it controls. For a contract the strategy is direct, since whoever deploys it can create owner addresses and have each of them approve it. The empirical findings do not rule this out. Section~\ref{sec:sybil-cost-model} formalizes the cost of attacking an allowance graph, where a closed farm of addresses that approve one another needs, unlike transfer wash trading, no token balance and pays little beyond gas. On the graph studied here such a farm buys a place in the upper part of the ranking for a few dozen approvals per window but would need more than a hundred thousand addresses to reach its top, because endorsement is concentrated: the highest-ranked cohort contract holds about 3\% of all stationary mass (Section~\ref{sub:sybil-bound}).

\dissertationSubheading[sec:practical-implications]{Practical implications and deployment scenarios}

Wallet providers, protocols and security tools that have to judge contracts can read either the flow into a contract or the authorizations it holds. The measured results support choosing the score by scenario over recommending a single one.

\dissertationSubsubheading[sub:deploy-transfer]{Transfer-hub and flow monitoring}

A protocol whose priority is describing current flow into contracts (for market-maker identification, routing analysis or transfer-based compliance heuristics) can read that flow from in-degree and in-value directly. AWP \parencite{do2023awp} adds propagation, time and value weights and restarts at active senders. It agrees with the transfer family in the same window (Table~\ref{tab:alignment-hybrid}), and out of window it has the highest coefficient of any score on later new senders, with a paired difference to the transfer in-degree above zero in both windows, although neither contrast was fixed in advance (Tables~\ref{tab:holdout-diff} and~\ref{tab:fresh-posthoc}). A protocol that wants to anticipate flow into a contract has reason to use AWP rather than the count.

\dissertationSubsubheading[sub:deploy-allowance]{Authorization-aware and low-latency refresh}

Where explicit spending authorization matters most (token vaults, allowance-gated integrations, or risk engines that treat approvals as acts of trust), protocols have reason to use allowance edges. The cheapest form is a spender's in-approve degree, which ranked new approval pairs better than EndorseRank in both windows, in contrasts not fixed in advance (Tables~\ref{tab:holdout-diff} and~\ref{tab:fresh-posthoc}), and needs no solver at all. EndorseRank is its propagated version. A refresh costs one solve over the allowance graph (Table~\ref{tab:benchmark-runtime}); the time comes from the size of that graph and would grow with any denser one. Its allowance alignment (Table~\ref{tab:alignment-hybrid}) and its holdout results (Table~\ref{tab:holdout-spenders}) make it usable where the operator wants a ranking that also reflects who the endorsing owners are. For most spenders, however, it reflects them only through how many spenders each owner approves and how recent and large each approval is (Section~\ref{sub:theory-closed-form}), weights not shown to improve any label over the count, and under uniform teleportation the propagated score is open to cheap inflation by closed farms (Section~\ref{sec:sybil-cost-model}). The variant with activity restarts did rank new approval pairs above the count, in comparisons made after the labels; an operator who wants a propagated allowance score has more reason to test that variant than the uniform one. In return for authorization semantics, operators accept a score that carries little transfer information.

\dissertationSubsubheading[sub:deploy-hybrid]{Evaluation design and combined use}

Whether to combine the two signals is a question about deployment and about method, and the evidence answers the two differently. An operator can run an allowance score for authorization-aware monitoring and AWP for flow, at the cost of both solves. Folding the two edge types into one walk did not pay on these data. C-PR walks both layers, so its solve costs more than AWP's (Section~\ref{sec:benchmark-results}), and although it ranks new approval pairs slightly better than EndorseRank in both windows, it trails AWP on new senders by a wider margin (Section~\ref{sub:finding-rq5}), so an operator free to choose has reason to keep the two scores apart. For an operator tied to a single transfer walk weighted by dollar amounts, rule B, applied to the revised scores, shows that adding allowance edges bought authorization information that held on new data, at a loss on inflow larger than the margin fixed in the registration. Because the pipeline is open, the analysis can be replicated on any ERC-20 chain that has allowance logs and a daily price series for its main tokens (Appendix~\ref{ch:appendix-eval}); integration with live risk oracles is left for future work (Chapter~\ref{ch:conclusion}).

\dissertationSubsubheading[sub:deploy-industry]{Industry adoption scenarios}

The scenarios below say for whom each score is useful and for whom it is not.

\begin{itemize}
\item Wallet software and approval prompts. An owner usually decides whether to trust a spender contract when its wallet shows the approval for signature, and weaknesses in how wallets present approvals make approval frauds easier \parencite{adamczyk2025approval}. A prompt could show, beside the amount requested, how many distinct owners hold a positive allowance to the spender, or its EndorseRank. That tells the owner how widely the contract is authorized. It tells nothing about what the contract will do with the allowance: a contract that collects approvals by phishing gains endorsement edges as an honest router does (Section~\ref{sec:approval-security}), and a new contract scores low whatever its code. A low score is a reason to look closer, and a high score is no reason to sign.
\item Protocol integrations. A protocol that routes user funds through external contracts, such as an aggregator choosing routers or a vault choosing strategies, has to decide which contracts to trust. The in-approve degree and EndorseRank rank how widely owners already authorize each candidate, and out of window they rank the contracts that go on to gain new approvals; AWP ranks those that go on to gain new senders (Table~\ref{tab:holdout-spenders}). Either is one input beside code review and audits, and neither says that a contract is safe to integrate.
\item Security tooling. Tools that list a user's open allowances and help revoke them could sort spenders by standing, so that the least endorsed are reviewed first. The revocation counts and the drain heuristic rise with every score on the spender holdout and track the volume of approvals a contract holds (Section~\ref{sec:holdout-results}), so no score here is a drainer or phishing detector. No score is a security audit either: a contract can contain malicious functions whatever its standing among owners (Section~\ref{sub:non-claims}).
\item Lending and margin protocols. Lending markets want to put more of their capital to use without cutting collateral \parencite{bartoletti2021lending}, and the stress of 2020 showed the cost of updating risk parameters late \parencite{gudgeon2020crisis}. When the borrower or trader is itself a contract, the screening question is whether it will be force-liquidated. On that outcome the comparison, its registered rules applied to the revised scores, found no count-level lead of allowance over transfer, although every trader contract that had received an approval avoided liquidation in both windows, in samples of 70 and 88 (Section~\ref{sub:finding-neutral}). A screen built on that pattern would need a registered test on a larger sample. The scores are not tools for predicting default.
\item Registries and incentive programs. Airdrop farming shows how readily on-chain rewards attract many-address strategies \parencite{luo2025airdrop}, and a program that lists or rewards contracts by their endorsements would face the same pressure. EndorseRank as defined here offers little protection below the top of its ranking, because a closed farm collects about three times its population share and can concentrate it on one target (Section~\ref{sec:sybil-cost-model}). Programs can require diverse incoming endorsements from owners with independent histories, discount mass that circulates inside closed clusters, or send teleportation and dangling mass only to a vetted seed set; none of these stops an attacker who buys approvals from honest owners.
\item Liquidity providers and market makers. To control inventory and spreads, these participants need to find the contracts that act as transfer hubs and see how flow is structured. Transfer counts and AWP \parencite{do2023awp} are the suitable tools for that goal, and their benefit is lower screening cost, not higher profit and loss.
\item Analytics and compliance providers. Exchanges, funds and compliance teams need contract rankings that outsiders can recompute. EndorseRank and AWP can be published for defined windows and parameters and audited, unlike proprietary machine learning models \parencite{packin2024credit}, provided that their Sybil exposure is disclosed and they are not presented as safety ratings or profit predictions.
\end{itemize}

\dissertationSubsubheading[sub:deploy-multiples]{Standardized scores as shared infrastructure}

A score shared by several parties is useful only if each of them can recompute it and see how it might be manipulated. EndorseRank and AWP can be recomputed from public logs, a versioned token list and a public daily price series with the open pipeline (Appendix~\ref{ch:appendix-eval}), and Sections~\ref{sec:sybil-cost-model} and~\ref{sec:regulatory-ethics} document how they can be manipulated.

\dissertationSubheading[sec:societal-impact]{Societal and practical impact}

In the DeFi market, vetting a counterparty is expensive, and the capacity to do it is unevenly spread. Wealthy actors may be able to pay for code audits and proprietary risk tools, while retail participants work from anecdotal evidence or skip vetting altogether. A score built from the permissions owners have already granted, and recomputable from public logs and public prices at a cost set by the size of the allowance graph (Section~\ref{sec:benchmark-results}), lowers the barrier to vetting the contracts that ask for those permissions. An approval that a user granted earlier is one route by which tokens are drained (Section~\ref{sec:approval-security}). A score centered on the spender contract therefore shows where that exposure lies, which transfer volume does not, although it does not say which spender will misuse an allowance (Section~\ref{sec:holdout-results}). The same score cuts the other way for developers: a new contract, however sound its code, starts with no endorsements, so a score used as a gate would favor contracts that are already widely approved.

For protocol developers, the results argue for replacing a single ``reputation'' construct with specific scores, each fitted to the decision it informs, which reduces the risk of using a flow score to decide authorization.

Where the decentralized credit scoring systems discussed by \textcite{packin2024credit} are opaque, a score built from public events and public prices, with a versioned token list and versioned parameters, supplies the explanatory information that regulators suggest should be required when DeFi protocols are analyzed \parencite{zetzsche2020defi}. It also brings a duty, described in Section~\ref{sec:regulatory-ethics}, to ensure that an ordinal, graph-based statistic is not mistaken for a security audit, a credit history, or an endorsement by whoever publishes it.

Any reputation metric that affects how value is allocated is open to exploitation; the next section models what it costs to manipulate EndorseRank.

\dissertationSubheading[sec:sybil-cost-model]{Sybil attack cost model on allowance graphs}

A Sybil attack exploits cheap pseudonyms to create a false impression of trustworthiness \parencite{douceur2002sybil}. On a transfer graph, a malicious actor can inflate in-degree and in-flow artificially by moving assets among accounts it controls, paying for this in transaction fees and possible inventory loss \parencite{cong2023washtrading}. Allowance graphs expose a different attack surface, since each edge records an authorization, not an actual transfer, and for a contract the attack is direct: whoever deploys it can create owner addresses and have each of them approve it. In this section we formulate a simple cost model for inflating EndorseRank with fake allowances. It builds on the optimal link farms of \textcite{gyongyi2005alliances} and on the result of \textcite{cheng2005sybilproof} that symmetric functions such as PageRank are not Sybil-proof (Section~\ref{sec:reputation-systems}), and adds the cost of an allowance edge and the size of the gain on this study's graph.

\dissertationSubsubheading[sub:sybil-setup]{Attack setup}

Let the attacker control a set of addresses $\mathcal{S}$, and write $\mathcal{H}=V\setminus\mathcal{S}$ for the honest addresses in the endorsement graph $G_{\text{approve}}=(V,E,w)$. An edge $(u,v)$ of weight $w_{uv}>0$ records that $u$ has authorized $v$ to spend its tokens. The weight is that of Equation~\eqref{eq:approve-weight}: for each of the five tokens, the time weight of the latest in-window approval times the bounded value weight of its US-dollar amount, summed over the tokens. With damping parameter $d\in(0,1)$, EndorseRank finds the PageRank vector $\mathbf{r}$ on $G_{\text{approve}}$ as the solution of \parencite{page1999pagerank}:
\begin{equation}
\label{eq:pagerank-sybil}
\mathbf{r} = d\, P^{\top} \mathbf{r} + \frac{1-d}{|V|}\mathbf{1},
\end{equation}
where $P$ is the transition matrix of Equation~\eqref{eq:transition}. A node with out-edges splits its row in proportion to its outgoing allowance weights. A dangling node, one without out-edges, has every entry of its row equal to $1/|V|$, so its mass is spread uniformly over all nodes \parencite{langville2006pagerank}. The fixed point is the same as that of the iteration in Equation~\eqref{eq:pagerank}, which the implementation runs and which feeds dangling mass back through the uniform teleportation vector. The attacker's objective is to maximize the mean EndorseRank of a target set $\mathcal{G}\subseteq\mathcal{S}$,
\begin{equation}
\label{eq:objective}
J(\mathcal{G}) = \frac{1}{|\mathcal{G}|}\sum_{i\in\mathcal{G}} r_i,
\end{equation}
subject to a budget constraint on the cost of the attack, which is defined next.

\dissertationSubsubheading[sub:sybil-cost]{Cost components}

The attacker creates a set of edges $E_{\mathcal{S}}$, all with sources in $\mathcal{S}$, and may also buy a set $E_{\mathcal{H}\to\mathcal{S}}$ of edges from honest owners into $\mathcal{S}$. Costs are counted over $q$ observation windows, the number of score refreshes for which the targets are meant to hold their rank. Gas is the only cost that every edge incurs. The other components arise only for particular edges or strategies.

\begin{enumerate}
\item Gas cost. Each \texttt{approve} or \texttt{permit} call consumes gas; write $c_{\mathrm{gas}}$ for the expected fee of one allowance update. EndorseRank reads only the \texttt{Approval} events inside its observation window, so an edge enters a window's graph only if its approval is emitted in that window. To keep $E_{\mathcal{S}}$ in place for $q$ windows, the attacker has to emit every approval again in each window, which also keeps its time weight near the maximum $\sigma(0)=0.86$, and the gas is paid $q$ times:
\begin{equation*}
C_{\mathrm{gas}} = q\,|E_{\mathcal{S}}|\,c_{\mathrm{gas}}.
\end{equation*}

\item Exposure cost. An approval locks no tokens, and the owner need not hold any balance of the token when it approves. The amount sets the weight but costs nothing: the value weight $V$ grows with the dollar amount, and an unlimited allowance takes its full value, so an owner that holds none of the token gives its edge the largest weight without exposing anything. After row normalization, moreover, a single out-edge gives $P_{uv}=1$ whatever its weight. The token list does not raise this cost, since approving one of the five tokens costs an attacker no more than approving any other token; what the list removes are events that a token contract can fabricate in the name of owners that never approved anything (Section~\ref{sec:amount-revision}), not approvals that a farm emits itself. An edge inside $\mathcal{S}$ therefore exposes nothing and costs nothing beyond gas. Exposure arises only when an attacker address that holds tokens approves a spender $v\notin\mathcal{S}$, for instance to look like an ordinary user of known contracts, since $v$ can then move those tokens. Such an edge also lets score mass leave $\mathcal{S}$. With $\ell_{uv}$ the expected cost to the attacker of one such exposure over one window,
\begin{equation*}
C_{\mathrm{exp}} = q\sum_{(u,v)\in E_{\mathcal{S}},\ v\notin\mathcal{S}} \ell_{uv},
\end{equation*}
which is zero for a closed farm.

\item Revocation and detection cost. Legitimate users revoke old allowances from time to time, and a farm whose approvals reappear unchanged in every window might be identifiable. Let $c_{\mathrm{rev}}$ be the expected cost per window of detection and of the revocations the attacker issues to avoid it. Over $q$ windows,
\begin{equation*}
C_{\mathrm{rev}} = q\,c_{\mathrm{rev}}.
\end{equation*}

\item Bought honest edges. The attacker can also pay honest owners to approve its addresses, for instance through an app that asks users for permission. A bought edge $(h,a)$, with $h\in\mathcal{H}$ and $a\in\mathcal{S}$, falls under the same window rule and costs, per window,
\begin{equation}
\label{eq:marginal}
c_{\mathrm{in}}(h,a) = c_{\mathrm{gas}} + p_h + \rho_h,
\end{equation}
where $p_h$ is the payment to the honest owner and $\rho_h$ is a risk premium the owner asks against the allowance being spent by the attacker.
\end{enumerate}

Over the $q$ windows, the total cost of the attack is
\begin{equation}
\label{eq:total-cost}
C = C_{\mathrm{gas}} + C_{\mathrm{exp}} + C_{\mathrm{rev}} + q\sum_{(h,a)\in E_{\mathcal{H}\to\mathcal{S}}} c_{\mathrm{in}}(h,a).
\end{equation}
For a closed farm, $C_{\mathrm{exp}}=0$ and $E_{\mathcal{H}\to\mathcal{S}}$ is empty, so $C=q\,(|E_{\mathcal{S}}|\,c_{\mathrm{gas}}+c_{\mathrm{rev}})$.

\dissertationSubsubheading[sub:sybil-bound]{Influence and a cost-effectiveness bound}

By the usual sensitivity results for PageRank, an edge moves $\mathbf{r}$ in proportion to the stationary mass at its source node and to its normalized outgoing weight \parencite{langville2006pagerank}. Consider first a closed farm: every node of $\mathcal{S}$ has at least one out-edge, all of these edges end in $\mathcal{S}$, and no honest owner approves an address in $\mathcal{S}$. Let $D$ be the set of dangling nodes, all of which then lie in $\mathcal{H}$, and let $r_D=\sum_{j\in D} r_j$ be their stationary mass. By Equation~\eqref{eq:pagerank-sybil}, a node $i\in\mathcal{S}$ receives $d\sum_{j\in\mathcal{S}} r_j P_{ji}$ from the farm, $(1-d)/|V|$ from teleportation and $d\,r_D/|V|$ from the dangling rows. The rows of the farm's nodes sum to one inside $\mathcal{S}$, so summing over $\mathcal{S}$ gives $\sum_{i\in\mathcal{S}} r_i = d\sum_{i\in\mathcal{S}} r_i + |\mathcal{S}|\,(1-d+d\,r_D)/|V|$, that is,
\begin{equation}
\label{eq:closed-cluster}
\sum_{i\in\mathcal{S}} r_i = \frac{|\mathcal{S}|}{|V|}\cdot\frac{1-d+d\,r_D}{1-d}.
\end{equation}
The first factor is the farm's population share, and the second says how many times that share the farm collects. The second factor equals one only when $r_D=0$, that is, on a graph without dangling nodes; otherwise the dangling rows pass part of the honest mass to the farm. Since the dangling nodes lie outside $\mathcal{S}$, $r_D\le 1-\sum_{i\in\mathcal{S}} r_i$, and substituting this into Equation~\eqref{eq:closed-cluster} gives a bound that holds on any graph:
\begin{equation}
\label{eq:cluster-mass}
\sum_{i\in\mathcal{S}} r_i \le \frac{|\mathcal{S}|}{(1-d)\,|V| + d\,|\mathcal{S}|}.
\end{equation}
For a farm that is small against $|V|$, the bound is close to $1/(1-d)$ times the population share. On the EndorseRank graph at $T_{\mathrm{obs}}$, with $6{,}219{,}063$ nodes and $11{,}210{,}114$ edges, the pipeline's own PageRank code gives $r_D=0.355$, $0.382$ and $0.404$ at $d=0.75$, $0.85$ and $0.95$, so a closed farm on this graph collects $2.06$, $3.17$ and $8.68$ times its population share. A direct check with the same code agrees: a ring of 20 addresses, each approving the next, added to the graph collects $3.17$ times its population share at $d=0.85$.

Equation~\eqref{eq:cluster-mass} limits the farm's total mass, not the score of one target, and the attacker can concentrate that mass. Take a star in which $f$ farm addresses each approve the target $t$ and $t$ approves each of them back. Every farm address $a$ has $t$ as its only out-neighbor, so $P_{at}=1$, and the row of $t$ sums to one over the farm. Write $c=(1-d+d\,r_D)/|V|$ for the mass that each node receives from teleportation and the dangling rows. The balance equations $r_a=c+d\,r_t P_{ta}$ and $r_t=c+d\sum_a r_a$ then give
\begin{equation*}
r_t = \frac{c\,(1+df)}{1-d^2}.
\end{equation*}
At $d=0.85$, $c=7.64\times10^{-8}$, the common score of every address that receives no approval (the constant $c$ of Section~\ref{sub:theory-closed-form}). With $f=10$ (eleven addresses and twenty approvals per window), the target reaches $r_t=2.62\times10^{-6}$ on the graph with the star added, the value of the closed form. That is above the median EndorseRank of the cohort contracts ($1.88\times10^{-7}$) and about half of their 90th percentile ($5.28\times10^{-6}$), but about one twelve-thousandth of the highest cohort score ($0.0318$). The score grows linearly in $f$: each further farm address adds $d\,c/(1-d^2)$ to $r_t$ and costs two approvals per window. Solved for $f$, the formula needs no farm at all for the median, since two addresses that approve each other already score $c/(1-d)$, above it; 22 farm addresses pass the 90th percentile, 2{,}074 the 99th ($4.86\times10^{-4}$), and about 135{,}895 reach the highest cohort score, before counting that a farm of that size would itself enlarge $|V|$ and lower $c$. Passing the median takes one pair of addresses and passing the 90th percentile a few dozen approvals per window, whereas reaching the top would take more than a quarter of a million approvals per window. The reason is concentration: the highest-ranked cohort contract holds about 3\% of the graph's stationary mass, more than four hundred thousand times $c$.

Honest in-edges are therefore not needed to reach the upper part of the ranking; without them, reaching its top takes a farm of the size just given, and bought honest edges add to the gain. When the farm has no out-edges to honest addresses, an edge $(h,a)$ bought from an honest owner $h$ adds, to first order,
\begin{equation*}
\Delta\sum_{i\in\mathcal{S}} r_i \approx \frac{d\,r_h\,P_{ha}}{1-d}
\end{equation*}
to the farm's mass, where $P_{ha}$ is the share of $h$'s row that the new edge takes. The gain depends on this share, and through $V$ on the dollar amount approved. A fresh approval of $z$ dollars bought from $h$ takes $\sigma(0)V(z)/(o(h)+\sigma(0)V(z))$ of $h$'s row, and an approval of a few hundred dollars or more, or an unlimited one, has $V(z)$ close to one and takes $\sigma(0)/(o(h)+\sigma(0))$. For the median owner in the graph that share is $0.957$, since the latest approvals of most owners carry small weights, through their age, their amount or both. The edge an attacker would buy is therefore a large or unlimited approval from an owner whose own approvals are old, small or few, preferably one with a high $r_h$. Such an approval puts the owner's balance of the token, up to the approved amount, at the attacker's disposal, so the risk premium $\rho_h$ in Equation~\eqref{eq:marginal} grows with what the owner holds; it can still be close to zero for an owner that holds none of the token. The allowance layer of C-PR weighs each allowance by its dollar amount up to the cap $U$, so there a bought approval takes a large share of an owner's row only if it is large against the owner's other allowances, and an approval that reaches the cap authorizes the spending of at least $U$ dollars of the owner's balance (Section~\ref{sub:coupled-graph}). The attack's cost-effectiveness is
\begin{equation}
\label{eq:efficiency}
\eta = \frac{\Delta J(\mathcal{G})}{C},
\end{equation}
with $\Delta J(\mathcal{G})$ the gain in the objective of Equation~\eqref{eq:objective} that the attack edges produce, held over the $q$ windows that $C$ in Equation~\eqref{eq:total-cost} pays for.

Equations~\eqref{eq:closed-cluster}--\eqref{eq:efficiency} show that EndorseRank, as defined here, is open to cheap inflation by closed farms. Under uniform teleportation and uniform redistribution of dangling mass, a closed cluster is not held to its population share. On this graph the inflation is cheap in the middle and upper part of the ranking and expensive at its top. The control that does bound closed clusters changes the walk itself: if teleportation and dangling redistribution go only to a vetted seed set, as in TrustRank \parencite{gyongyi2004trustrank}, a cluster outside that set gains mass only through edges from honest owners, which the attacker has to buy at the price of Equation~\eqref{eq:marginal}. AWP's restart rule gives the transfer side no protection either: a farm whose addresses send small amounts to one another raises their activeness and so draws restart mass to itself. The variant of EndorseRank with activity restarts, which ranked new approval pairs above the count (Section~\ref{sub:finding-rq4}), inherits the same exposure, since every farm address that approves raises its own approving activity. The Sybil-adjusted stability proxies (Section~\ref{sub:proxy-sybil}) flag repeated transfers from a few addresses but not a farm of fresh addresses, so they cannot replace an adversarial assessment, which we leave to future work.

\dissertationSubheading[sec:regulatory-ethics]{Regulatory and ethical considerations}

On-chain reputation metrics raise questions of transparency as well as equity. Legal work has shown that decentralized finance markets remain bound by disclosure and consumer protection rules even though settlement happens on-chain \parencite{zetzsche2020defi}. \textcite{packin2024credit} add a warning of their own: decentralized credit scoring may perpetuate black-box harms when models cannot be interpreted or contested, or when they are linked to identity databases without due process. Both EndorseRank and AWP are computed deterministically from on-chain events and a public daily price series, with open-source code and a versioned token list, so anyone can recompute them (Appendix~\ref{ch:appendix-eval}). That openness partly answers the black-box criticism. It leaves the distributional concern untouched, because contracts with few approvals, every new contract among them, score lower whatever their code.

Ethically, these scores should never be presented as security ratings, as audits or as guarantees that a contract is safe to approve, nor as creditworthiness measures in the traditional sense. They are rank-based indicators of a contract's position in the defined graph. If a wallet, protocol or registry acts on score thresholds, for instance by warning users or by declining an integration, it should document the graph definitions behind the scores, the window and the limitations, above all the exposure to Sybil attacks and the fact that a phishing contract collects approvals as an honest one does. A score that wrongly marks a contract as little endorsed can harm its developers, and one that wrongly marks it as widely endorsed can harm the owners who trust it. A regulatory framework centered on interpretability would prefer this kind of ranking to proprietary machine learning models like those reviewed by \textcite{kotb2023credit}, as long as providers make no exaggerated claims about what it predicts \parencite{packin2024credit}.

Privacy on an open blockchain is limited by the ledger's very nature \parencite{narayanan2016bitcoin}, and computing the scores draws on nothing beyond public data. The scored objects are contracts. The externally owned accounts that approve, pay or are paid by them enter the graphs only as owners, senders and recipients; none of them is scored, ranked or reported, and no attempt is made to link any address to a real-world person (Section~\ref{sec:ethics}). Owners' approvals still enter the computation, and the score of a contract that one person controls could say something about that person, so the scores should not be treated as permanent, fingerprint-like identifiers.

\dissertationSubheading[sec:threats-validity]{Threats to validity}

The findings reported here should be read against the limitations below, which are grouped under construct, internal, external and adversarial validity.

\dissertationSubsubheading[sub:threat-construct]{Construct validity}

Seven points concern what the scores and their checks measure.

\begin{itemize}
\item These scores are not credit ratings or security ratings: no default labels from uncollateralized lending and no exploit or audit labels enter the analysis.
\item Same-window checks rest on local allowance and transfer statistics of the same contracts.
\item Holdout construct. A future new approval is still an endorsement event and a future new sender a flow event: both lie outside the score window, yet neither is an independent measure of trust or safety. Being counts of arrivals, they also favor any score that tracks the current count of counterparties (Section~\ref{sub:finding-rq4}).
\item Unit of analysis. EndorseRank scores the receiving side of an approval, and every scored address here is a contract: the matched cohort holds the contracts that received approvals of one of the five tokens from at least three owners, the spender cohorts the contracts that held a positive allowance of one of them at the freeze, and the trader cohorts the GMX~V2 contract accounts with at least three closes (Section~\ref{sub:cohorts}). The same-window results say nothing about contracts that fewer owners approve or that owners approve only on other tokens, and no result says how the scores would order externally owned accounts. Claims about trading outcomes rest on the two small trader cohorts (Section~\ref{sub:finding-neutral}).
\item Account type. An address counts as a contract when it holds code at one block read on 9~October~2026, after both label windows (Section~\ref{sub:cohorts}). A contract could once remove its own code with \texttt{SELFDESTRUCT} \parencite{wood2025yellowpaper}; since EIP-6780 \parencite{ballet2023eip6780}, only a contract created and destroyed in one transaction loses its code. A contract that destroyed itself before that change has no code at the block read, counts as an externally owned account and is left out of every cohort, so if such contracts differ from the ones that survive, for instance because some of them were scams, the cohorts lean toward contracts that still exist. In the other direction, an address that received approvals before code was deployed at it counts as a contract for the whole window, and an account with an EIP-7702 delegation \parencite{buterin2024eip7702} counts as an externally owned account although it runs code.
\item Own-construct agreement is partly mechanical. EndorseRank smooths spender in-approve degree over the whole graph, and AWP does the same for incoming transfers weighted by their value, so each score is expected by construction to agree with its own family. The contrasts of Table~\ref{tab:tau-diff} were computed for this reason; with few ties in either score they compare coefficients with similar attainable maxima, but they remain checks of the instruments.
\item Value weights. EndorseRank and AWP pass each US-dollar amount through $V(z)$ with the slope $b=\ln 3/m$, where $m$ is the median dollar amount of the transfers or of the finite latest allowances up to $t_1$, rounded to two significant digits (Section~\ref{sec:amount-revision}). An amount at the median weighs one half and one at twice the median $0.8$, so the weights record the value of the approvals and transfers an edge carries as well as their number and age (Table~\ref{tab:alignment-transfer}). The slopes are choices of the revision (Section~\ref{sub:threat-internal}). In the registered analysis, with $b=1$ on raw base units, $V$ was within $10^{-4}$ of one for almost every amount, so the weights there recorded number and age and almost nothing about value (Appendix~\ref{sec:map-registered}). The layers of C-PR add dollar amounts without the bound of $V$, capping each allowance at $U=20{,}000$ dollars, so an unlimited allowance weighs as much as a large finite one and an owner that holds both kinds splits its row among them in proportion to their capped amounts (Section~\ref{sub:coupled-graph}). On both spender holdouts EndorseRank has a slightly higher coefficient on both labels than the allowance layer walked alone (Tables~\ref{tab:holdout-spenders}, \ref{tab:fresh-holdout} and~\ref{tab:fresh-posthoc}). No paired contrast between the two was computed, so the gap suggests, without showing, that how approvals are weighted matters for this construct. The comparison mixes the value weight with the time weight, which only EndorseRank applies to approvals, and neither weighting is shown to be the right one. Such weighting choices propagate into $\mathbf{r}$ and into the rank correlations \parencite{langville2006pagerank}.
\end{itemize}

\dissertationSubsubheading[sub:threat-internal]{Internal validity}

Six points concern whether the comparisons isolate what they are read as showing.

\begin{itemize}
\item Shared observation window. Every method and proxy is computed over 1~October~2023 to 30~June~2026, with the time decay anchored at the end. Market regimes and protocol changes peculiar to those months may push alignment up or down. The checks on damping and on the value slopes measure sensitivity to these parameters, but none of them varies the calendar window, and the two holdout freezes, three months apart, vary it only a little.
\item Cohort selection. The matched cohort is every contract that received non-zero approvals of one of the five tokens from at least three owners in the window. It is a subset of the plan's cohort, taken from that cohort's approval logs and code check, and was fixed in the revision before any revised score was computed; results may not carry over to contracts with fewer approvals or with approvals of other tokens only. Not every cohort contract is in both graphs (Section~\ref{sub:cohorts}). Keeping the missing contracts as isolated nodes instead of scoring them zero changes no coefficient (Table~\ref{tab:tie-sensitivity}), and AWP's zero for spenders without a transfer edge does not explain its holdout results (Table~\ref{tab:holdout-receiving}). The sample-definition stages show how strongly the coefficients react when the evaluation set changes (Table~\ref{tab:robustness-sample-size}), which is why they are not read as robustness evidence.
\item Baseline fidelity. AWP follows the published method in its edge weights and its restarts (Section~\ref{sub:awp-graph}). The paper gives no parameter values, so $k$, $t_0$ and $b$ are this study's choices, $b$ set in the revision (see value weights above). The paper scored Ethereum transactions; here AWP scores transfers of five ERC-20 tokens on Arbitrum One, valued in US dollars, within the ego network of the cohort. EndorseRank takes AWP's weights, with its own slope for allowances, and restarts uniformly, as fixed in the analysis plan, and the variant with AWP's restarts is reported beside it (Table~\ref{tab:endorserank-restarts}).
\item Order of analysis. The cohorts, scores, labels, statistics and rule A were fixed in the analysis plan of 9~October~2026, before any log was extracted. Rule B with its margin, its further results F4--F6, the two sensitivity analyses and the comparison on trader outcomes were registered at 16:49~UTC on the same day, after the W0 results were known and before the July--September logs were extracted (Section~\ref{sub:fresh-holdout}). These dates order the registered analysis; on the revised scores every rule was applied after the registered results on the same labels were known (next point). Every other comparison was computed after the relevant labels were known: in W0, the degree contrasts of the allowance layer, the S-PR contrasts, the comparisons of C-PR with its transfer layer on new approval pairs and with AWP, and the restriction to spenders that had received a transfer; in both windows, on the spender labels, the degree contrasts of EndorseRank and AWP, the comparisons of C-PR with EndorseRank, the comparisons of EndorseRank with AWP and every contrast of the activity-restart variant. They are reported with intervals, and none of them decides a question.
\item Revision of the amounts. The amounts were revised on 10~October~2026, after the registered results of the same window, W0 and W1 were known (Section~\ref{sec:amount-revision}), so the revised analysis is post hoc: its choices, the five tokens, the price source, the slopes $b$ and the cap $U$, were made with those results in view, and a revision made after the results can favor one outcome over another. Three facts limit that risk without removing it. The revision was committed to the public repository before any revised score was computed; it derives the slopes and the cap from quantiles of the events up to $t_1$; and the registered analysis reached the same verdicts under rules A, B and R1--R3 (Appendix~\ref{sec:map-registered}, Table~\ref{tab:registered-vs-revised}). The slopes and the cap remain choices of the revision: dividing or multiplying the slopes by ten changes the same-window results little (Table~\ref{tab:robustness-slope}), and no other cap was tried.
\item Timing environment. All runtimes are wall-clock times from one machine (Section~\ref{sec:benchmark-results}), and their standard deviations cover only repeated solves of one configuration. Variation between machines is not measured, so the comparison is stated as a time ratio beside the edge ratio, with $|V|$ and $|E|$ of each graph.
\end{itemize}

\dissertationSubsubheading[sub:threat-external]{External validity}

Five points limit how far the findings carry beyond this sample.

\begin{itemize}
\item Single chain and window. The evidence comes from Arbitrum One alone, over 33 months of scored activity and two three-month label windows. Extending the findings to other rollups, Layer-1 chains or lending-centered venues would need the same evaluation on those venues, since DeFi venues differ in protocol design \parencite{werner2022sok} and in market structure \parencite{schar2021defi}.
\item Tokens and prices. The graphs hold the approvals and transfers of five tokens only, WETH, USDC, USDT0, ARB and WBTC (Table~\ref{tab:usd-tokens}). Every other token is left out, legitimate ones included, so a contract that owners approve only on such tokens is outside the cohort, and native ether, which moves without \texttt{Transfer} logs, enters the graphs only as WETH. Every price comes from one public source, the daily series of the DefiLlama API \parencite{defillamanddocs}. On the dates sampled for the check it agreed closely with the Chainlink feeds read on the chain \parencite{chainlinknddatafeeds}, but an error in the series on another date would pass into every amount of that day (Appendix~\ref{sec:data-tokens}).
\item Ego network. Each graph keeps only the approvals and transfers of the five tokens in which a cohort contract takes part (Section~\ref{sub:cohorts}). An owner's approvals of contracts outside the cohort, and transfers between two addresses outside it, are missing, so owners' out-strengths, senders' activeness and the paths the walks can take are those of the ego network and not of the chain. Paths of more than one step exist only through cohort contracts that both approve and are approved (Section~\ref{sub:theory-closed-form}), and the findings about propagation do not carry over to graphs with longer endorsement paths.
\item Scaling stages. The scaling stages are subsets of the cohort and its ego network, so they show how the solve time follows $|E|$ up to the size of that graph and not beyond (Section~\ref{sec:benchmark-results}). Same-window alignment remains anchored to the matched cohort and the holdout to spenders.
\item GMX contract traders. Few contract accounts close positions on GMX~V2 in the label windows, and the trader cohorts hold 70 contracts in W1 and 88 in W0, so every trader result has wide intervals. The comparison on outcomes built from neither edge type says nothing about the externally owned accounts that also trade on GMX~V2, which this study does not score.
\end{itemize}

\dissertationSubsubheading[sub:threat-adversarial]{Adversarial validity}

Three points concern deliberate manipulation and deception.

\begin{itemize}
\item Sybil and wash trading. Rankings may be manipulated through adversarial allowance or transfer patterns \parencite{douceur2002sybil}; the Sybil-adjusted proxies only partly address this, and the cost model of Section~\ref{sec:sybil-cost-model} is analytical, with one ring and one star added to the observed graph, not an empirical red-team evaluation.
\item Phishing. A contract that obtains approvals by deception gains endorsement edges as an honest one does (Section~\ref{sec:approval-security}), and nothing in the scores separates the two.
\item Strategic approval behavior. Once a score is live, the developers of a contract may arrange approvals to it to farm rank. An observational $\tau$ estimated before deployment may overstate how well the ranking resists such behavior.
\end{itemize}

\dissertationSubheading[sec:limitations]{Limitations}

On this sample, EndorseRank's strengths are agreement with the allowance checks, a ranking distinct from AWP's and, on spender contracts, a ranking of new approval pairs in both windows that the scores built on transfers alone do not match. Its lower cost follows from a graph with fewer than half of the transfer graph's edges and is a property of that graph, not of the method. Its weaknesses are a ranking that carries little transfer information, a holdout coefficient below that of the raw in-approve degree it smooths in both windows, exposure to closed farms below the top of its ranking, and nothing to say about contracts that few owners approve or that owners approve only on other tokens. These strengths and weaknesses are those of the revised analysis, whose tokens, prices, slopes and cap were chosen after the registered results were known (Section~\ref{sub:threat-internal}); the registered analysis reached the same verdicts on every rule.

The coupled operator was tried at three fixed interior values of $\lambda$, with uniform teleportation and a per-node convex mixture at the transition level, on two layers weighted by dollar amounts, with each allowance capped at $U$. Layers weighted as in EndorseRank and AWP, node-dependent or learned weights, mixtures of the stationary vectors instead of the transition rows, and other Markov chains went untested, so both answers to RQ5 apply only to that family of methods on this dataset. Deployment remains bound by the considerations of Section~\ref{sec:regulatory-ethics}.

EndorseRank draws on no off-chain evidence about who controls an address, and neither does any other score here; the only off-chain input is the daily price series that values the amounts. Attestation tooling such as Human Passport uses identity checks, biometrics and vouching to decide whether an address belongs to a distinct person \parencite{humanpassportnddocs}. That question has no direct counterpart for a contract, and EndorseRank gives no guarantee that the owners who approve a contract are controlled by different actors, so it should be read as an ordering of contracts by the approvals they hold and not as a Sybil defense. Counting only the approvals of owners that pass such a gate before computing EndorseRank is a plausible combination, but it is untested here: the gate's own weaknesses, such as credential farming \parencite{bordeianu2026sybil}, would carry over to it, and evaluating the combination needs labeled Sybil data that this study did not collect.

These limitations narrow what can be concluded. For contracts that at least three owners approved on WETH, USDC, USDT0, ARB or WBTC on Arbitrum One in these months, EndorseRank is an allowance-specific reputation score for spender contracts. It is not a security audit or a credit score, it does not resist Sybil farms, and it does not replace protocol-level collateralization \parencite{gudgeon2020crisis}. The allowance edge it rests on is worth adding to on-chain reputation, and the approval count is the cheapest way to use it. Chapter~\ref{ch:conclusion} draws the conclusions and names the work that would extend these bounds.
